\section{近期研究（2026 年 8--11 月）}
\paragraph{Stipple：重新利用双目惯导前端的关键帧做增量泼溅}
2026.08~\cite{2608_00931}。Stipple 以 Basalt 跟踪并产出位姿、关键帧和闭环。当前地图先渲染关键帧，系统据尚未覆盖的像素决定何时添加原语，并把昂贵的双目深度计算限于部分锚点帧。高斯训练与双目惯导跟踪并行。

作者在 EuRoC 与 Monado 的公开序列和相应双目基线上报告渲染质量与吞吐量。高帧率的比较应区分前端跟踪速度与增量建图频率。此外，这条路线依赖双目标定和惯导输入。

\paragraph{EvTrajGS：将无位姿事件流的相机估计连续化}
2026.08~\cite{2608_08585}。离散事件姿态独立优化容易在时间上出现不一致。EvTrajGS 从通用管线得到粗位姿，用连续时间函数表示相机轨迹；渲染时汇总邻近时刻的轨迹证据，联合更新相机与高斯。重加权采样让训练关注监督不足的区间。

在由 Replica 五个场景合成的事件上，论文相对各指标第二名报告平均 PSNR 增益 3.84 dB。真实 TUM-VIE 部分给出轨迹误差与渲染定性图，不能把 Replica 的渲染数字转用于真实事件。方法仍使用粗轨迹初始化，不能理解为完全没有相机估计输入。

\paragraph{EndoMD-SLAM：退化观测下暂停静态内镜地图的更新}
2026.08~\cite{2608_08949}。水滴和污物可能遮住内镜镜头，暂态假象也可能被回填进长期地图。EndoMD-SLAM 用观测可靠性控制分支选择。退化帧触发历史干净关键帧检索和重定位，并暂停静态地图写入。自监督的静态与暂态分离有助于把持续解剖结构与短时光学伪影分开。

论文从 C3VDv2 选出十段含镜头水或碎屑标签的序列比较跟踪和重建。结果说明了这类镜头遮挡下的表现，重定位仍需可匹配的历史关键帧，不能直接推广到长期没有干净视线的操作。

\paragraph{在线高斯建图中的透射率保持、尺度初值和残差增密}
2026.08~\cite{2608_14902}。在线增密的每个阶段迭代有限，直接复制原语可能累积不透明度偏差。作者在 Photo-SLAM 的建图端加入保持合成透射率的增密，按深度和相机内参估计单像素的米制足迹来初始化尺度，并将新原语优先放入可从深度回投的高残差位置。

实验同时测量地图渲染和相机跟踪，并逐项添加模块比较。它沿用分离的跟踪架构，主要改动在建图侧。论文按 Photo-SLAM 工具的全轨迹评估图像，不能把数值与只评估训练视图的表格直接并列。

\paragraph{MotionGS-SLAM：事件流调节曝光积分中的高斯投影}
2026.08~\cite{2608_15024}。MotionGS-SLAM 以事件流估计局部运动，让投影高斯沿运动方向拉伸，同时随整体运动增加曝光内位姿的采样密度。两类调节进入渲染侧，用来解释观测帧的运动模糊；这个处理无需把底层三维高斯全部拉长来拟合拖影。

在合成 EventReplica 上，论文比较可用清晰真值的渲染质量和轨迹误差；真实测试包括三段 Color-DAVIS346 暗光采集。复用该流程需要与帧同步的事件数据，合成图像指标不能代替更多真实场景的定量验证。

\paragraph{AquaFlow：在水下单目流中同步估计位姿、地图与介质}
2026.08~\cite{2608_22906}。AquaFlow 从水下单目流进入在线高斯建图。用水下数据适配的几何模型提供位姿跟踪与深度初始化线索，同时将可见颜色随距离变化拆给衰减和后向散射模型，使场景本征外观和介质现象不必都储存为视角相关颜色。

论文组织了 62 段视频，但定量重建只是其中可评测的子集，且作者报告 SeaThru-NeRF 上低于离线 UW-GS。应分开比较实时与离线约束，也不能把整个测试套件看作均有米制轨迹和参考几何。

\paragraph{CGS-SLAM：让单目惯导局部高斯子图在服务器对齐}
2026.08~\cite{2608_26868}。各客户端从单目与 IMU 生成局部视觉惯导轨迹，用单目米制深度辅助高斯子图尺度。协作阶段共享关键帧编码以发现重叠，再在服务器用 VGGT 帮助子图配准，合成全局表示。

作者报告笔记本跟踪建图约每秒十帧，原实验将较重的 Depth Pro 放在独立服务器，关键帧通信量取决于采样节奏，合并地图也需要传输。端侧运行的结果与完全离线或不依赖服务器的系统应分别理解。

\paragraph{SCOUT-SLAM：在动态场景中分别估计跟踪和建图的不确定性}
2026.09~\cite{2609_14634}。动态场景中，错误位姿会损坏地图，损坏的地图又会误导下一次跟踪。SCOUT-SLAM 从各关键帧提取 DINOv2 特征，由共享网络估计建图不确定性，再用叠加在共享网络上的低秩适配器估计跟踪不确定性，按像素重新加权位姿优化残差。两项估计通过不同监督更新，共用图像特征，同时避免直接把当前地图误差当成跟踪可靠度。

论文在 TUM RGB-D、Bonn Dynamic 和 Wild-SLAM MoCap 上比较轨迹或建图效果。TUM 的静态渲染指标排除了用 SAM3 标记的人和椅子区域，不能解读为对整幅动态场景的重建。附录另在 DROID-W 评估，其中 Downtown 1 和 4 两条超过 100 米的轨迹因显存不足未报告。

\paragraph{PanoGS-SLAM：在球面域优化全景光度与相机轨迹}
2026.09~\cite{2609_17387}。PanoGS-SLAM 将一张全景与统一三维高斯地图匹配来估计相机位姿，并在球面条件下增量写入新观测。等距矩形投影在高纬有采样面积失衡，光度损失据此调整贡献，地图则由深度先验引导初始化。

作者在合成全景 SynPano 及实采宽视场 PALVIO 对比定位，并用受控视场实验分析宽角覆盖。结果需结合输入的球面相机几何；同名全景设备未必覆盖相同的竖直角度。

\paragraph{VGGT-GS SLAM：联合细化未标定单目子图的几何与内参}
2026.09~\cite{2609_19628}。VGGT-GS SLAM 先用前馈几何给每个时间窗初始化高斯子图，再用可微束调整共同优化高斯、相机外参、子图共享内参和畸变。其后在全局位姿图中使用高斯子图间的相似变换约束，降低各窗独立初始化导致的尺度与漂移不一致。

论文在 TUM-RGBD、7-Scenes、ScanNet 等室内视频上评价定位，输入仅用 RGB。作者指出性能受几何先验质量与标定可观测性影响，窗口也带来一定延迟。

\paragraph{EliGSiR：按地图当前需要动态分配 RGB-D 在线预算}
2026.09~\cite{2609_20348}。EliGSiR 对持续 RGB-D 建图中的视图、分辨率和原语增长分别分配预算。地图变化后可重新排序历史视图，输入负载升高时可降低监督分辨率再恢复，在重复观测证实几何缺失的位置才追加容量。

它在 Replica、TUM RGB-D 等序列评估留出视图，并把输入流内、有限清尾和额外后优化的状态明确区分。多数受控比较供应同一参考位姿，因此绘图误差不能解释为实时定位精度，跨不同试验组也不宜直接比较绝对 PSNR。

\paragraph{RawSLAM：用线性高动态范围观测进行高斯跟踪与建图}
2026.09~\cite{2609_20589}。RawSLAM 将相机 RAW 转成线性 HDR 输入，用对数方式参数化高斯颜色，并调整光度梯度的亮度权重，使跟踪和高斯训练能保留高光与阴影的信息。这与将八位色调映射图像直接作为观测的做法不同。

作者在十段室内序列中比较原生 HDR 与常规图像处理，报告改善随场景照明而变化。PSNR、SSIM 和 LPIPS 在色调映射后的图像上计算，因此这些指标不能当作直接的线性辐亮度误差。

\paragraph{Cube-Splat：将全景图像分面后共用一个位姿变量}
2026.09~\cite{2609_21347}。直接将全景图像送入针孔高斯渲染器会受到投影模型限制。Cube-Splat 将全景帧分解成固定朝向、共享光心的虚拟透视视图，并复用透视高斯渲染。每个分面产生的梯度用伴随变换映射到同一个位姿变量，因而不会把各分面误当成可以独立运动的相机。

作者在 SynPano、OmniBlender 和 PALVIO 共 31 个序列中报告轨迹误差和视图质量，并给出逐序列结果。该表示依赖准确的全景标定和固定分面朝向，拼接误差会影响接缝附近的光度一致性。系统针对大体静态的环境，未融合惯性或深度观测，缺纹理与重复结构下仍可能漂移。

\paragraph{Elevator-VIGS：把电梯轿厢运动与轿厢内相机运动分开}
2026.09~\cite{2609_23491}。相机在轿厢内看到的是相对于轿厢的运动，而惯导响应包含轿厢整体运动。Elevator-VIGS 在束调整中加入沿重力方向的运输状态，通过视觉语言和深度线索识别电梯事件，使轨迹与高斯地图能在移动坐标系下更新。

作者测试十六段真实电梯记录，并用模拟电梯序列补充过程真值。真实轨迹主要在上下电梯前后的控制点评价竖向误差，中间全程轨迹的证据来自仿真。设计针对沿重力方向运动的电梯。

\paragraph{BayesianGS-SLAM：把预测不确定性用于跟踪、建图和选帧}
2026.09~\cite{2609_24140}。BayesianGS-SLAM 保留 LoopSplat 的 RGB-D 高斯地图和在线更新流程，将颜色及深度的传感器噪声与不透明度引起的地图表示不确定性合并为预测量。跟踪按预测不确定性调整残差，选帧则看当前观测相对地图预测有多意外，从而减少对已解释视角的重复更新。

原文报告 TUM 的平均绝对轨迹误差为 3.12 厘米，ScanNet 为 6.7 厘米；三条 TUM 序列上的插入关键帧和建图调用均约减少七成。渲染指标在训练关键帧上计算，只反映这些视图上的拟合效果。机器人采集部分为定性演示，不能代替大规模独立轨迹验证。

\paragraph{双协方差高斯 SLAM：让渲染与配准共享中心而各自建模}
2026.09~\cite{2609_25746}。光度优化得到的高斯协方差表示适合渲染的形状，不一定反映深度测量的可信度。该方法为共享中心的高斯保存两套协方差，一套按光度优化，一套表示 RGB-D 测量不确定性。后者用于深度的广义 ICP 残差，也给锚定到附近高斯的图像角点提供置信度，补充几何较弱时的跟踪线索。

在 TUM RGB-D、ScanNet 和 Replica 中，作者分别统计绝对轨迹误差、相对姿态误差与渲染指标，并报告困难手持或室外例子。系统缺少大场景长期一致性所需的全局束调整，且深度传感误差随距离增大，因此不能将近距离的跟踪收益外推到远景主导的环境。

\paragraph{ArborSplat：以几何与类别预算维护果园语义地图}
2026.09~\cite{2609_26315}。ArborSplat 用 LiDAR 里程计提供相机位姿，在线维护带类别分数的高斯。图像掩码结合双目点和行间地面几何形成语义监督，并在预算受限时按类别分配部分新原语。高斯地图不反馈到跟踪端。

论文在苹果与梨的三个阶段比较在线地图。薄树干在训练视图上可识别，但留出视图的树干 IoU 在完整路线中仍低于 0.07，这一差距限制了新视角的细结构语义稳定性。

\paragraph{HyperLoop：用负结果界定高光谱闭环检测的实际收益}
2026.10~\cite{2610_19742}。HyperLoop 在基于高斯的稠密 SLAM 回环模块中并列使用 RGB 描述子与光谱描述子，以严格同意或加权机制筛选闭环。它改变的是地点识别和回环决策，高斯场的建图观测没有直接改用高光谱。

在 MMOTS-WS 上，最佳融合的对齐 ATE RMSE 是 5.04 厘米，而单独使用 SIFT-VLAD 为 4.95 厘米。作者把相对 NetVLAD 的轨迹收益归因于 RGB 描述子替换。这个负结果针对 600 至 860 纳米与当前场地，不能据此断言其他波段无效。

\paragraph{GOFusion：在占据子图中锚定高斯并积累信息量}
2026.10~\cite{2610_25261}。GOFusion 用视觉惯导估计位姿和占据子图，再以体素与表面法向引导活动子图的高斯初始化。占据层支持规划空间查询，高斯层保留外观。光栅化过程中累计的对角 Fisher 近似用作原语裁剪和选视的信息指标。

论文在 TUM RGB-D、uHumans2 和真实飞行记录中分别评价。飞行演示是给定平台和传感器的实例，Fisher 近似用于相对信息量排序，并不是导航安全保证。

\paragraph{SG²SLAM：用流式点图和三维关联连接户外高斯子图}
2026.10~\cite{2610_36312}。SG²SLAM 让冻结的前馈模型按连续帧窗口预测局部位姿、点图与置信度，用三维先验同时服务跟踪和高斯建图。高斯核受视线条件调节，窗口间通过全局相似变换图纠正单目尺度和累计漂移。

论文在 KITTI 和 Waymo 等户外序列用 Sim(3) 对齐后的轨迹误差评价，不能把对齐误差解释为原生绝对尺度验证。作者也把雨夜等非结构化场景的稳定性留作后续研究。

\paragraph{LingGS：把流式前馈模型当作虚拟 RGB-D 传感器}
2026.10~\cite{2610_36853}。LingGS 让冻结的流式模型预测内参、深度和图像特征，然后以显式对应、重投影和深度误差估计相机运动。高斯后端消费连续图像、深度和这些位姿，不直接采用前馈模型相机头输出的轨迹。

作者在 7-Scenes、Oxford Spires 等室内外数据上比较定位和渲染，也检验不同流式主干下的显式运动求解。该设计把位姿误差与几何预测解耦，但仍依赖深度、匹配特征及相机内参的可用性。

\paragraph{DPGS-SLAM：区分真实观测与暂存的生成背景}
2026.11~\cite{2611_12038}。DPGS-SLAM 用光流差异和 SAM2 找动态区域，由真实静态区域支持相机跟踪。异步扩散修补可为遮挡背景生成暂存高斯，但这些原语带来源记录，不参与位姿与观测地图优化；有效真实静态观测到达后可替换它们。

作者用 BONN、TUM 评价轨迹，并用 StaDy4D 检验静态背景。不可见区域仍可能长期保留错误外观或几何，异步补全也不保证实时。保留来源的意义在于让真实测量和生成内容承担不同角色。

\paragraph{CoSphere-SLAM：球面闭环与多机器人全景子图融合}
2026.11~\cite{2611_22235}。CoSphere-SLAM 把每台机器人的全景转到立方体视面用于局部估计，并在服务器按球面方向匹配不同朝向的关键帧。全局相似变换一致性用于过滤对称外观导致的伪闭环，同时协调各局部地图的尺度与位姿，再融合高斯。

论文比较多个室内和跨采集会话全景数据，在表中分别列出有有效轨迹的场景统计和全部场景统计。比较时应读清失败处理与是否多机器人联合运行，不能只拿成功子集的轨迹均值反映整体稳定性。
